% =============================================================================
% Unaffordable NILs, 2015-2024: share above the 40% threshold and ranking (2024)
% Generated from data/work/aff_nil_year.dta (data after commit f4d35f2 on main).
% Sample: NILs with OMI rent observed in every year 2015-2024 (75 NILs), the
% period with a consistent OMI zoning. Standalone: compile with pdflatex.
% To include in the paper, copy the text between \begin{document} and \end{document}.
% =============================================================================
\documentclass[11pt,a4paper]{article}
\usepackage[utf8]{inputenc}
\usepackage[T1]{fontenc}
\usepackage{lmodern}
\usepackage[margin=2.5cm]{geometry}
\usepackage{amsmath}
\usepackage{booktabs}
\usepackage{tabularx}
\usepackage{array}
\newcolumntype{L}{>{\raggedright\arraybackslash}X}
\begin{document}

\subsection*{How many NILs are out of reach?}

A NIL is classified as unaffordable when the annual OMI rent of a standard
60~m$^2$ dwelling exceeds 40\% of the median net income of Milanese taxpayers
($D_{it}=1$). We restrict this part of the analysis to 2015--2024, the period
in which the OMI zoning is consistent: the zones redrawn in 2014 remain in
force until 2024, and the 2024 split of zones D15 and D30 is mapped back to
the same boundaries (Section~2.2). The sample contains the 75 NILs with a rent
observed in every year of the period.

Table~\ref{tab:share} reports, for each year, the share of NILs above the
threshold and the share of residents living in them. In 2015, 21.3\% of NILs (16 of 75) were above the threshold, hosting 24.7\% of residents; in 2024 the figures are 36.0\% (27 of 75) and 43.2\%. Over the decade, 16 NILs are above the threshold in every year, 48 are never above it, and 11 cross it at least once.

Table~\ref{tab:rank} lists the five most and the five least affordable NILs
in 2024. With a median net income of EUR~21,559, the threshold corresponds to a rent of about EUR~11.98 per m$^2$ per month. The most affordable NILs are peripheral, in the south
(Gratosoglio), the north-west (Quarto Oggiaro) and the west (Muggiano, Quinto
Romano, Baggio): there a 60~m$^2$ flat absorbs between a quarter and a third
of the median net income. The least affordable are central: in Duomo and Brera
the same flat would absorb more than 80\%. Because income is a single
city-wide value, the ranking is a ranking of OMI rents, and NILs in the same
OMI zone receive similar values.

\begin{table}[htbp]
\centering
\caption{Share of NILs and of residents above the 40\% threshold, 2015--2024}\label{tab:share}
\small
\begin{tabular}{@{}lrr@{}}
\toprule
Year & NILs above threshold (\%) & Residents in those NILs (\%) \\
\midrule
2015 & 21.3 & 24.7 \\
2016 & 25.3 & 29.5 \\
2017 & 25.3 & 29.5 \\
2018 & 22.7 & 27.1 \\
2019 & 25.3 & 29.5 \\
2020 & 26.7 & 29.9 \\
2021 & 24.0 & 28.2 \\
2022 & 26.7 & 29.7 \\
2023 & 30.7 & 34.9 \\
2024 & 36.0 & 43.2 \\
\bottomrule
\end{tabular}

\vspace{4pt}
\begin{minipage}{0.8\linewidth}
\footnotesize\emph{Notes:} 75 NILs with OMI rent observed in every year
2015--2024. Threshold: $A_{it}=12\cdot60\cdot R_{it}/Y^{net}_t>0.40$.
Residents: DS205 counts for the same year. Source: \texttt{aff\_nil\_year.dta}
(module 18).
\end{minipage}
\end{table}

\begin{table}[htbp]
\centering
\caption{The five most and the five least affordable NILs, 2024}\label{tab:rank}
\small
\begin{tabularx}{\linewidth}{@{}rLrrrrr@{}}
\toprule
Rank & NIL (code) & Residents & Rent & Annual rent & $A_{i,2024}$ & $A_{i,2015}$ \\
 & & 2024 & (EUR/m$^2$/month) & 60 m$^2$ (EUR) & (\%) & (\%) \\
\midrule
\multicolumn{7}{@{}l}{\emph{Most affordable}} \\
1 & Gratosoglio - Q.re Missaglia - Q.re Terrazze (41) & 18,912 & 7.69 & 5,534 & 25.7 & 24.3 \\
2 & Quarto Oggiaro - Vialba - Musocco (76) & 32,744 & 7.73 & 5,569 & 25.8 & 22.8 \\
3 & Muggiano (54) & 3,045 & 9.32 & 6,714 & 31.1 & 24.9 \\
4 & Quinto Romano (62) & 4,744 & 9.33 & 6,714 & 31.1 & 24.9 \\
5 & Baggio - Q.re degli Olmi - Q.re Valsesia (55) & 30,860 & 9.68 & 6,971 & 32.3 & 26.3 \\
\midrule
\multicolumn{7}{@{}l}{\emph{Least affordable}} \\
75 & Duomo (1) & 16,226 & 25.11 & 18,082 & 83.9 & 89.4 \\
74 & Brera (2) & 18,361 & 24.53 & 17,659 & 81.9 & 79.7 \\
73 & Guastalla (4) & 15,470 & 21.05 & 15,158 & 70.3 & 70.4 \\
72 & Porta Garibaldi - Porta Nuova (9) & 5,817 & 18.73 & 13,488 & 62.6 & 62.1 \\
71 & Magenta - S.Vittore (7) & 18,026 & 17.76 & 12,789 & 59.3 & 58.7 \\
\bottomrule
\end{tabularx}

\vspace{4pt}
\begin{minipage}{\linewidth}
\footnotesize\emph{Notes:} rank among the 75 NILs (1 = most affordable,
75 = least affordable). Rent: OMI midpoint for civil dwellings in normal
condition, area-weighted from OMI zones to NILs, mean of the semesters of 2024.
$A_{it}$: annual rent of 60~m$^2$ over the median net income of Milanese
taxpayers (EUR~21,559 in 2024; EUR~17,370 in 2015). All ten NILs are on the same
side of the 40\% threshold in 2015 and 2024. Ranks 3 and 4 (Muggiano, Quinto
Romano) have virtually identical rents. For Porta Garibaldi--Porta Nuova (9),
2015 rents of zone C14 are imputed from zone C12, from which C14 was carved
(quoted only from 2018S2). Source: \texttt{aff\_nil\_year.dta} (module 18).
\end{minipage}
\end{table}

\end{document}
